\section{Experimental Setup}
Primary quantitative results in this chapter come from the reported public-corpus experiments and
the shared Android adaptive-window comparison. The evidence package includes dataset and protocol
manifests, model reports, privacy ablation, metadata leakage, encrypted-flow fingerprinting,
privacy-student and federated prototype reports, performance gates, and the app-selectable model
comparison.

The core evaluation principle is to avoid optimistic random-row validation. Model and privacy
reports use grouped splits keyed by dataset source, environment, session, app family, and time
bucket wherever the report type supports them. The separate \texttt{evaluation-protocol.json} audit
covers six grouped split strategies over a sampled 250,000-window subset to verify that the
methodology itself is source-aware.

\section{Datasets and Trace Collection}
The corpus is a public-only merge of seven normalized datasets: Android Mischief, Android Spyware,
CICAndMal2017, ITC-Net-Blend-60 scenario E, a Parrot-derived metadata slice, SDNCampus, and
Westermo. The merged corpus contains 6,034,487 flows and yields 904,371 60-second windows for
feature extraction. \Cref{tab:dataset-sources} summarizes the dataset composition reported by
\texttt{dataset-manifest.json}.

\begin{table}[htbp]
  \centering
  \caption{Public-only corpus composition.}
  \label{tab:dataset-sources}
  \small
  \begin{tabular}{p{4.1cm}rrp{5.1cm}}
    \toprule
    Dataset source & Flow rows & Sessions & Label profile \\
    \midrule
    Android Mischief & 3,011 & 8 & anomaly-only slice \\
    Android Spyware & 4,924 & 1 & mixed benign/anomalous \\
    CICAndMal2017 & 2,587,130 & 2,119 & mixed benign/anomalous \\
    ITC-Net-Blend-60 E & 17,199 & 211 & benign-only \\
    Parrot metadata slice & 373,566 & 1 & benign-only \\
    SDNCampus & 3,000,000 & 30 & benign-only \\
    Westermo & 48,657 & 3 & mixed benign/anomalous \\
    \bottomrule
  \end{tabular}
\end{table}

The corpus spans malware-heavy Android traffic, benign campus/application traffic, and industrial
network traces. It is therefore more heterogeneous than a single-device lab trace. This makes the
benchmark harder, but closer to the cross-source generalization problem that MANTA is intended to
study.

\section{Metrics}
\begin{itemize}[leftmargin=*]
  \item \textbf{Detection metrics:} precision, recall, F1, ROC-AUC, PR-AUC, and confusion matrix.
  \item \textbf{Privacy/utility metrics:} anomaly utility under \texttt{off}, \texttt{low},
  \texttt{medium}, and \texttt{strict} views.
  \item \textbf{Metadata leakage metrics:} strongest-attacker accuracy and normalized leakage for
  app ID, app family, dataset source, context bucket, and destination behavior.
  \item \textbf{Observer leakage metrics:} encrypted-flow sequence fingerprinting accuracy and
  normalized leakage.
  \item \textbf{Performance metrics:} single-window and batch feature-extraction latency.
\end{itemize}

Normalized leakage is reported on a \(0\)--\(1\) scale after accounting for the baseline difficulty
of the inference task. A value near \(0\) means that the attacker is near the configured baseline or
chance level. A value near \(1\) means that the representation is highly identifying for the target
attribute. The metric is an empirical privacy-risk score, not a formal privacy guarantee.

\section{Android-Deployable Experiment Methodology}
The Android-deployable experiments test how far the local Android path can improve while using only
features that the endpoint can compute at runtime. They reuse the public corpus, generate
Android-style adaptive windows, enforce an explicit Android feature contract, and evaluate
deployable random-forest model files that can be interpreted by the Android JSON tree scorer.

The methodology added five safeguards after the early local results proved too weak:
\begin{itemize}[leftmargin=*]
  \item \textbf{Feature-contract audit:} every added feature must come from Android packet/flow
  metadata, destination insight, or recent local windows. Server-only and label-only fields are
  excluded to avoid train/deploy skew.
  \item \textbf{Window-construction alignment:} training windows are generated to match the runtime
  app-window builder more closely, including adaptive windows for sparse/service traffic.
  \item \textbf{Split and balance audits:} reports track source, family, and label imbalance so
  global metrics cannot hide a single dominant dataset or app family.
  \item \textbf{Label-conflict handling:} rounded near-duplicate windows with both benign and
  malicious labels are downweighted or excluded rather than treated as clean contradictory
  supervision.
  \item \textbf{Window and episode metrics:} per-window precision/recall/F1 are paired with
  episode-level alert metrics, because the Android app surfaces correlated alert episodes rather
  than isolated CSV rows.
\end{itemize}

The Android UI uses descriptive names for the deployable models. ``High-recall RF'' is the
recall-oriented local model. ``Low-FPR RF'' is the conservative low-false-positive model. ``Adaptive
hybrid RF'' is the runtime combination of both. These names describe operating behavior; the exact
versioned files and commands are retained in the reproducibility appendix.

\section{Anomaly-Detection Results}
\Cref{tab:model-matrix} summarizes the headline anomaly-detection variants under one shared Android
adaptive-window comparison. The backend was retrained on the same Android-style adaptive feature
cache and updated to consume the full portable Android feature vector rather than the older partial
remote feature payload. Its threshold is calibrated on the Android validation holdout with an
explicit false-positive budget instead of an earlier max-F1 or recall-heavy threshold. The
calibrated boosted-tree backend reaches precision \(0.898\), recall \(0.878\), F1 \(0.888\),
PR-AUC \(0.962\), and ROC-AUC \(0.966\) on the balanced holdout, with false-positive rate reduced
to \(0.100\).

\begin{table}[htbp]
  \centering
  \caption{Representative model results under the shared Android adaptive-window holdout.}
  \label{tab:model-matrix}
  \small
  \begin{tabular}{>{\raggedright\arraybackslash}p{4.0cm}cccc}
    \toprule
    Model & Precision & Recall & F1 & PR-AUC / ROC-AUC \\
    \midrule
    High-recall RF & 0.916 & 0.933 & 0.925 & 0.972 / 0.974 \\
    Privacy RF & 0.888 & 0.934 & 0.910 & 0.898 / 0.948 \\
    Adaptive hybrid RF & 0.952 & 0.845 & 0.895 & 0.977 / 0.979 \\
    Backend boosted tree, calibrated & 0.898 & 0.878 & 0.888 & 0.962 / 0.966 \\
    Low-FPR RF & 0.967 & 0.707 & 0.817 & 0.973 / 0.976 \\
    \bottomrule
  \end{tabular}
\end{table}

The results show usable signal in metadata-only windows. The strongest balanced F1 comes from the
Android High-recall RF path. Adaptive hybrid RF sits between the recall-oriented and quiet local
models: it recovers much of the Low-FPR RF recall while keeping a lower false-positive rate than the
High-recall RF path. The backend tree model also remains strong after the feature-contract mismatch
is removed. Threshold calibration eliminates the previous pathological \(1.000\)-recall operating
point. The backend is still noisier than the Low-FPR RF path, but it is no longer the worst
false-positive option.

\section{Experiment Progression and Negative Results}
The experimental sequence was intentionally iterative. The first Android-local experiments used a
small endpoint-safe feature set and simple deployable models. They proved that the app could score
traffic locally, but their recall and F1 were too weak for an interactive security tool. The
TFLite-compatible one-class path was also retained as a deployable interface experiment, but its
early result had low recall. This showed that a mobile inference format is not enough. The feature
contract, learning objective, and score calibration also have to match the operational detection
problem.

Several later attempts were documented but not selected:
\begin{itemize}[leftmargin=*]
  \item \textbf{Larger expanded-window caches:} 500,000-window and early 250,000-window adaptive
  cache runs exceeded practical memory limits because the builder materialized too many derived
  columns across the raw flow table. The memory-safe row builder completed, but required a longer
  build path for the expanded experiment cache.
  \item \textbf{Boosted-tree export:} the deployable boosted-tree candidate completed, but did not
  beat the random-forest variants on the selected split. It was therefore documented as a useful
  comparison rather than used as the default local model.
  \item \textbf{HistGradientBoosting benchmark:} the non-deployable histogram gradient boosting
  model estimated the value of a stronger tree learner, but it was slightly worse than the selected
  teacher-weighted random forest and lacked an implemented Android export path.
  \item \textbf{Pure service specialist:} a service-only specialist improved precision but lost too
  much recall. The result suggested that service behavior was internally broad and label-noisy. A
  specialist without better labels or alert-episode logic became overly conservative.
  \item \textbf{Full hybrid backend and hybrid-model sweeps:} the Android runtime hybrid was
  implemented, but the Python exported-tree evaluator timed out while repeatedly scoring both
  90-tree forests and sweeping episode policies. Full backend adaptive-window retraining with the
  heavier hybrid dual-channel family exceeded practical runtime in this environment. The lighter
  boosted-tree backend completed on the same Android adaptive-window cache. These failures are
  evaluation/training-path bottlenecks rather than mobile-runtime blockers.
\end{itemize}

These negative results shaped the final interpretation. The largest gains came from fixing
train/deploy mismatch, adding Android-computable temporal and destination context, auditing noisy
labels, and evaluating alert episodes. Changing the classifier family alone was less important than
aligning the data representation with the Android runtime and the user-visible alert policy.

\section{Android-Deployable Model Development}
The Android deployable path was iterated because early on-device models were not acceptable for an
interactive app. The final reported app comparison therefore focuses on the app-selectable models in
\Cref{tab:current-app-comparison}, while this section explains the development logic qualitatively.

High-recall RF expanded the Android feature contract from the original compact set to local
transport, destination-role, destination-transition, and per-app baseline-deviation features. It
also added source/family balancing, hard-example mining, teacher-weighted random forests, and
label-conflict downweighting. The main failure mode was not a single missing classifier family.
Service windows with similar volume and diversity statistics were being treated as
indistinguishable. Adding local context gave the model more axes for separating repeated service
behavior from anomalous shifts.

Low-FPR RF then added short temporal memory, destination-intelligence features, a malware
specialist, stricter service threshold policy, and stronger ambiguous-label handling. Its purpose
was not to maximize global F1. Its purpose was to reduce the user-visible false alerts that a purely
recall-oriented model could create.

Adaptive hybrid RF was added as a runtime hybrid rather than as a new trained model. It includes the
recall-oriented and quiet models. Malware-like and remote-access traffic keeps the more sensitive
score, while service/system/other-app traffic is damped toward the quieter score. In the shared
offline comparison this yields precision \(0.952\), recall \(0.845\), F1 \(0.895\), and FPR
\(0.043\), showing the intended compromise between sensitivity and alert-noise control.

Episode-level analysis motivated the Android runtime policy that stores weak pending evidence and
requires persistence, evidence diversity, or a high-confidence bypass before surfacing an alert.
This matters because user-visible alerts are correlated episodes over time, not isolated CSV rows.

The Android runtime was also changed after the offline model iterations. Explicit local model modes
can surface single-window validation candidates instead of waiting for repeated evidence. Known-danger destination evidence can bypass the normal model threshold. The statistical detector has a
cold-start risk floor for unusually suspicious first observations. These changes address a separate
problem from model F1: user-visible alerts can be suppressed by stability gates even when the model
score is nonzero. They are implementation safeguards and require a dedicated on-device alert study
before being treated as offline detection-quality results.

\subsection{Shared-Protocol App Model Comparison}
Because the later Android models were trained and calibrated at different times, the final
comparison uses one shared Android adaptive-window cache, one source-aware holdout split, and each
model's exported threshold. \Cref{tab:current-app-comparison} reports the
balanced-holdout view. This is more informative for comparing precision and false-positive rate than
the natural holdout because the natural sampled holdout is anomaly-heavy.

\begin{table}[htbp]
  \centering
  \caption{App-selectable models under the shared Android adaptive-window methodology.}
  \label{tab:current-app-comparison}
  \small
  \begin{tabular}{lccccc}
    \toprule
    Model & Precision & Recall & F1 & FPR & Threshold \\
    \midrule
    High-recall RF & 0.916 & 0.933 & 0.925 & 0.085 & 0.825 \\
    Privacy RF & 0.888 & 0.934 & 0.910 & 0.117 & 0.995 \\
    Adaptive hybrid RF & 0.952 & 0.845 & 0.895 & 0.043 & 0.790 \\
    Backend boosted tree & 0.898 & 0.878 & 0.888 & 0.100 & 0.999 \\
    Low-FPR RF & 0.967 & 0.707 & 0.817 & 0.024 & 0.790 \\
    \bottomrule
  \end{tabular}
\end{table}

This comparison clarifies the operating trade-offs. High-recall RF is the strongest reported single
model by balanced F1. Privacy RF clears the target precision/recall gate under the same methodology,
but it is not better than every normal model. Adaptive hybrid RF improves recall over the quiet path
while keeping false positives well below the recall-oriented path. Low-FPR RF is conservative: it
produces fewer false positives but misses more positive windows.

\section{Release-Privacy Utility Trade-off}
The release-privacy ablation results are shown in \Cref{tab:ablation}. In the reported merged-corpus
report, \texttt{off} and \texttt{low} retain the full Android feature contract and therefore have
identical utility. \texttt{medium} and \texttt{strict} reduce exact destination and fine-grained
feature information. For release leakage, these reduced views lower exact app re-identification,
but they also reduce detection utility. This is the expected privacy-utility trade-off once the
report is aligned with the final feature set.

\begin{table}[htbp]
  \centering
  \caption{Anomaly utility under privacy-reduced views.}
  \label{tab:ablation}
  \begin{tabular}{lcccc}
    \toprule
    View & Precision & Recall & F1 & PR-AUC / ROC-AUC \\
    \midrule
    \texttt{off} & 0.651 & 0.962 & 0.776 & 0.759 / 0.857 \\
    \texttt{low} & 0.651 & 0.962 & 0.776 & 0.759 / 0.857 \\
    \texttt{medium} & 0.565 & 0.871 & 0.685 & 0.661 / 0.756 \\
    \texttt{strict} & 0.515 & 0.893 & 0.653 & 0.582 / 0.703 \\
    \bottomrule
  \end{tabular}
\end{table}

\subsection{Hashing Versus Deleting Reduced Features}
An additional experiment compared three privacy-reduction strategies under the same Android
adaptive-window cache and the same source-aware holdout split. The goal was to distinguish semantic
deletion from pseudonymization. In the \emph{deleted} representation, the retained medium features
are bucketed and every suppressed Android feature is set to zero. In the \emph{hashed}
representation, those same suppressed features are kept only as deterministic salted hash buckets.
The \texttt{medium-plus} representation instead uses broad bucketization over many Android
features while zeroing direct identity/threat-tag ratio fields.

\Cref{tab:hash-delete} shows the balanced-holdout result. Hashing suppressed features recovers
substantial utility compared with deletion, but it also increases semantic leakage. This supports a
key privacy interpretation: hashing is pseudonymization, not removal. It can preserve stable
fingerprints that help detection, but those same fingerprints can help an attacker infer app family
or app identity.

\begin{table}[htbp]
  \centering
  \caption{Balanced-holdout comparison of reduced-feature deletion and hashing.}
  \label{tab:hash-delete}
  \small
  \begin{tabular}{lccccc}
    \toprule
    Representation & Precision & Recall & F1 & App-family acc. & App-ID acc. \\
    \midrule
    Full reference & 0.974 & 0.954 & 0.964 & 0.557 & 0.344 \\
    Medium-plus bucketed & 0.978 & 0.937 & 0.957 & 0.585 & 0.280 \\
    Hashed suppressed features & 0.971 & 0.868 & 0.917 & 0.590 & 0.239 \\
    Deleted suppressed features & 0.934 & 0.707 & 0.805 & 0.471 & 0.208 \\
    \bottomrule
  \end{tabular}
\end{table}

\section{Metadata Leakage from Released Views}
\Cref{tab:leakage} summarizes the strongest metadata leakage tasks for the main release-privacy
views. The result is mixed. Exact app re-identification drops substantially from the \texttt{off} view to
\texttt{medium} and \texttt{strict}. However, app-family, destination-behavior, and dataset-source
leakage remain high even in the reduced views. This is why the release-privacy gate for
\texttt{medium} and \texttt{strict} still fails.

\begin{table}[htbp]
  \centering
  \caption{Strongest metadata leakage by privacy view (normalized leakage).}
  \label{tab:leakage}
  \small
  \begin{tabular}{lcccc}
    \toprule
    View & App ID & App family & Dataset source & Destination behavior \\
    \midrule
    \texttt{off} & 0.282 & 0.722 & 0.927 & 0.999 \\
    \texttt{low} & 0.272 & 0.799 & 0.902 & 0.988 \\
    \texttt{medium} & 0.151 & 0.678 & 0.871 & 0.450 \\
    \texttt{strict} & 0.135 & 0.603 & 0.646 & 0.580 \\
    \bottomrule
  \end{tabular}
\end{table}

In both \texttt{medium} and \texttt{strict}, exact app re-identification falls below the configured
threshold. The app-family and destination-behavior checks remain false. The reduced views therefore
improve release privacy for one target while leaving broader semantic leakage unresolved.

\section{Local and Remote Models Under Privacy Views}
The local/remote privacy utility report separates two questions. The first is the real Android
runtime path: the local privacy model receives the normal Android feature window and applies its own
internal \texttt{medium-plus} input transform before scoring. The second is a release-view path: a
remote consumer receives an already reduced \texttt{medium} or \texttt{strict} export view and tries
to reuse a full-feature scorer. These are not equivalent.

On a balanced 5,000-window utility sample, the app-facing local privacy runtime path reaches
precision \(0.904\), recall \(0.956\), and F1 \(0.930\). The backend scorer reaches precision
\(0.724\), recall \(0.987\), and F1 \(0.835\) on the same balanced sample, but with a much higher
false-positive rate. Under already reduced \texttt{medium}/\texttt{strict} release views, the direct
full-feature local and remote scorers still have no useful operating mode because most expected raw
inputs have been removed. Reduced export views therefore require a dedicated student/reduced-view
model, an internally transformed local scorer, or explicit threshold/policy calibration rather than
naive reuse of the full-feature production scorer.

\section{Observer Leakage from Encrypted Traffic}
The observer benchmark evaluates what a passive on-path observer could infer from the encrypted
traffic sequence. This benchmark is deliberately separate from the release views, because the
observer sees the original traffic before any export-time coarsening. \Cref{tab:observer} shows that
observer leakage remains severe.

\begin{table}[htbp]
  \centering
  \caption{Observer leakage from encrypted-flow sequence fingerprinting.}
  \label{tab:observer}
  \begin{tabular}{lcccc}
    \toprule
    Task & Accuracy & Normalized leakage & Top-5 accuracy & Strongest model \\
    \midrule
    App ID & 0.544 & 0.544 & 0.719 & MLP \\
    App family & 0.810 & 0.783 & 0.979 & MLP \\
    Dataset source & 0.984 & 0.981 & 1.000 & MLP \\
    Context bucket & 0.046 & 0.044 & 0.107 & MLP \\
    \bottomrule
  \end{tabular}
\end{table}

The reduced privacy views improve the representation that MANTA exports. They do not protect
against an observer who sees timing, burst structure, and sequence-level metadata on the network
path. The experiment marks the boundary between feature-level privacy controls and transport-level
traffic-analysis defenses.

\section{Privacy Student, Federated Path, and Performance}
The privacy-student model improves representation privacy only partially. In the prototype
\texttt{medium}-view report, adversary accuracy is reduced strongly for app ID but remains notable
for app family, dataset source, and destination behavior. The federated path remains a measured
collaboration prototype rather than the main detector in the evaluated configuration.

Performance gating is also incomplete. \Cref{tab:gates} summarizes the thesis-relevant gate
outcomes. Single-window extraction meets the p95 target (\(22.5\) ms) but misses the p99 target
(\(87.9\) ms versus the desired \(40\) ms). The feature-matrix construction itself is fast, but
end-to-end single-window build time is not yet at the intended deployment threshold.

\begin{table}[htbp]
  \centering
  \caption{Selected thesis gates and measured outcomes.}
  \label{tab:gates}
  \scriptsize
  \begin{tabular}{@{}p{3.0cm}p{2.4cm}p{6.3cm}p{1.2cm}@{}}
    \toprule
    Gate & Target & Measured outcome & Result \\
    \midrule
    Latency p95 & \(\leq 40\) ms & \(22.5\) ms & Pass \\
    Latency p99 & \(\leq 40\) ms & \(87.9\) ms & Fail \\
    Normalized app-ID release leakage & below threshold & \texttt{medium}: \(0.151\); \texttt{strict}: \(0.135\) & Pass \\
    Semantic release leakage & below threshold & app-family and destination leakage remain high & Fail \\
    Observer-leakage audit & low leakage & app ID \(0.544\); dataset source \(0.981\) & Fail \\
    Detector quality & deployment-grade & F1: \(0.925\), \(0.910\), \(0.895\), \(0.888\), and \(0.817\) for High-recall RF, Privacy RF, Adaptive hybrid RF, Backend, and Low-FPR RF; backend FPR \(0.100\) & Mixed \\
    \bottomrule
  \end{tabular}
\end{table}

\section{Interpretation}
The evaluation supports several findings and rejects a stronger privacy interpretation.
\begin{itemize}[leftmargin=*]
  \item Metadata-only anomaly detection is technically feasible on a large public corpus.
  \item The calibrated backend run shows that backend inference remains a strong centralized option
  when remote scoring is acceptable.
  \item The Android-deployable comparison shows that the local path improves substantially when
  training windows, runtime windows, feature contracts, and label-conflict handling are aligned. The
  engineering conclusion is therefore ``local detection is sensitive to deployability alignment'',
  not simply ``local detection is too weak.''
  \item Release-privacy transformations can reduce exact app re-identification while preserving
  reasonable utility.
  \item Observer leakage remains severe without transport-level defenses.
  \item Runtime alert quality is not identical to per-window model quality. Episode grouping,
  evidence persistence, known-danger bypasses, and false-positive budgets change what the user
  actually sees.
  \item MANTA should be treated as an evaluation framework and research prototype rather than as a
  product-ready privacy-preserving IDS.
\end{itemize}
